\documentclass[11pt]{article}
\usepackage{amsmath, amssymb, amsthm}
\usepackage{geometry}
\usepackage{booktabs}
\usepackage{hyperref}
\geometry{margin=1in}

\newtheorem{theorem}{Theorem}
\newtheorem{lemma}{Lemma}
\newtheorem{definition}{Definition}
\newtheorem{proposition}{Proposition}

\title{Mathematical Supplement: TLMOS Hybrid Kernel\\
\large Formal Derivation of the Modified Pressure Scheduler}
\author{Traffic Light Management Operating System (TLMOS)}
\date{}

\begin{document}
\maketitle

\section{Notation and System Model}

Let an intersection be modeled as a directed graph $G = (\mathcal{L}, \mathcal{M})$
where $\mathcal{L}$ is the set of approach \emph{lanes} and $\mathcal{M}$ is the set
of \emph{movements} (directed pairs $(u, d)$ with $u, d \in \mathcal{L}$).

\begin{definition}[Lane State]
Each lane $u \in \mathcal{L}$ at time $t$ is characterised by:
\begin{itemize}
  \item $q_u(t) \in \mathbb{Z}_{\geq 0}$ — queue length (vehicles)
  \item $W_u(t) \in \mathbb{R}_{\geq 0}$ — head-of-queue waiting time (seconds)
  \item $\mu_u \in \mathbb{R}_{>0}$ — saturation flow rate (vehicles/second)
\end{itemize}
\end{definition}

\begin{definition}[Signal Phase]
A phase $s \in \mathcal{S}$ is a maximal set of non-conflicting movements:
\[
  s \subseteq \mathcal{M}, \quad \forall (u,d),(u',d') \in s: \text{movements do not conflict}
\]
Let $\mathcal{S}_n$ denote the set of all admissible phases for intersection $n$.
\end{definition}

\section{The Hybrid Kernel Scheduling Problem}

The kernel solves the following constrained optimisation at each control interval
$\Delta t = 1\,\text{s}$:

\begin{equation}
  \boxed{s^*(t) = \arg\max_{s \in \mathcal{S}_n} P_s(t)}
  \label{eq:argmax}
\end{equation}

where the \emph{modified pressure} $P_s(t)$ is:

\begin{equation}
  P_s(t) = \underbrace{\alpha \sum_{(u,d) \in s} \mu_{u,d}
  \bigl(q_u(t) - q_d(t)\bigr)}_{\text{Max-Pressure (throughput) term}}
  + \underbrace{\beta \max_{u \in s} W_u(t)}_{\text{WFQ (fairness) term}}
  \label{eq:pressure}
\end{equation}

\textbf{Hyperparameters:}
\begin{itemize}
  \item $\alpha \in [0,1]$ — weight on throughput maximisation
  \item $\beta \in [0,1]$ — weight on starvation prevention; $\alpha + \beta = 1$
\end{itemize}

\paragraph{OS Analogy.}
Equation~\eqref{eq:pressure} is the traffic analogue of the
\emph{Virtual Finish Time} used in Weighted Fair Queuing in network packet
schedulers \cite{demers1989analysis}. The term $\beta \cdot W_u(t)$ assigns a
priority bonus proportional to how long the head-of-queue vehicle has waited,
exactly as WFQ assigns earlier virtual finish times to flows with long-pending packets.

\section{Deadlock Formalisation}

\subsection{Coffman Condition Mapping}

Traffic gridlock satisfies all four Coffman \cite{coffman1971system} necessary
conditions for deadlock:

\begin{table}[h]
\centering
\begin{tabular}{lll}
\toprule
\textbf{OS Condition} & \textbf{Traffic Equivalent} & \textbf{TLMOS Response} \\
\midrule
Mutual Exclusion   & Conflicting streams cannot co-occupy & Phase conflict matrix \\
Hold \& Wait       & Vehicle occupies space, waits for next & Queue-state monitoring \\
No Preemption      & Cannot teleport vehicle out of lane & Flush Phase (rerouting) \\
Circular Wait      & $A \to B \to C \to A$ intersection chain & WFG cycle detection \\
\bottomrule
\end{tabular}
\caption{Coffman conditions mapped to urban gridlock.}
\end{table}

\subsection{Wait-for Graph}

\begin{definition}[Wait-for Graph]
The WFG at time $t$ is a directed graph $\mathcal{W}(t) = (\mathcal{L}, E(t))$ where
$(u, v) \in E(t)$ iff:
\[
  q_u(t) > 0 \;\land\; d_u = v \;\land\; q_v(t) > 0
\]
Lane $u$ is blocked by lane $v$ when $u$ has queued vehicles whose downstream
destination $d_u = v$ is itself congested.
\end{definition}

\begin{theorem}[Gridlock Equivalence]
A deadlock (gridlock) exists at intersection $n$ at time $t$ if and only if
$\mathcal{W}(t)$ contains a directed cycle.
\end{theorem}

\begin{proof}
($\Rightarrow$) If gridlock exists, vehicles in a cycle $\ell_1 \to \ell_2 \to
\cdots \to \ell_k \to \ell_1$ each hold their lane and wait for the next, directly
yielding a cycle in $\mathcal{W}(t)$.

($\Leftarrow$) A cycle $(\ell_1, \ell_2, \ldots, \ell_k, \ell_1)$ in $\mathcal{W}(t)$
means $q_{\ell_i}(t) > 0$ for all $i$, and each $\ell_i$ is blocked by $\ell_{i+1}$.
Under the No Preemption condition (vehicles cannot be removed), no lane can drain
while the cycle persists, satisfying the OS definition of deadlock. \qquad $\square$
\end{proof}

\subsection{Flush Phase Resolution}

When the kernel detects a cycle $C = (\ell_1, \ldots, \ell_k)$:

\begin{enumerate}
  \item \textbf{Context Save}: Snapshot all queue states and phase timers.
  \item \textbf{Bottleneck Identification}: Select the victim lane
        $\ell^* = \arg\min_{\ell \in C} q_\ell(t)$ (minimum queue $\Rightarrow$
        fastest to drain).
  \item \textbf{Force-Green}: Override the scheduler and grant green to the
        phase serving $\ell^*$ for $\min\_green$ seconds.
  \item \textbf{Context Restore}: Return to normal $P_s(t)$ scheduling.
\end{enumerate}

This is analogous to an OS \emph{aborting the victim process} with the lowest cost
(minimum work lost) to release a resource and break the deadlock cycle.

\section{Bounded Delay Guarantee}

\begin{theorem}[Anti-Starvation for $\beta > 0$]
For any lane $u$ with a non-zero arrival rate $\lambda_u > 0$, if $\beta > 0$,
then the maximum waiting time $W_u(t)$ is bounded:
\[
  \exists\; T_{\max} < \infty : \forall\, t,\; W_u(t) \leq T_{\max}
\]
\end{theorem}

\begin{proof}[Proof Sketch]
Suppose for contradiction that $W_u(t) \to \infty$. Since $\beta > 0$, the
WFQ term $\beta \cdot W_u(t)$ in $P_s(t)$ grows without bound for every phase $s$
serving $u$. Therefore there exists time $T$ such that $P_{s_u}(T)$ exceeds the
pressure of any competing phase with bounded wait times. By~\eqref{eq:argmax}, the
scheduler must select $s_u$ at or before $T + \max\_green$, contradicting the
assumption. $\square$
\end{proof}

\noindent\textbf{Interpretation.} Pure Max-Pressure ($\beta = 0$) is throughput-optimal
\cite{tassiulas1992stability} but provides no fairness guarantee. The TLMOS kernel
with $\beta > 0$ achieves a Pareto-optimal balance: throughput is near-maximised
while bounded delay is guaranteed for all movements—a contribution not present in
existing MP literature.

\section{Hyperparameter Sensitivity}

The tradeoff surface over $(\alpha, \beta)$ with $\alpha + \beta = 1$ can be
characterised empirically. Based on simulation results (see
\texttt{metrics\_output/results.csv}):

\begin{itemize}
  \item $\beta = 0$ (pure MP): Optimal throughput, unbounded delay on minor streets.
  \item $\beta = 0.3$: Near-optimal throughput, Jain fairness index $\approx 0.91$.
  \item $\beta = 1.0$ (pure WFQ): Fair but suboptimal throughput under saturation.
\end{itemize}

We recommend $(\alpha, \beta) = (0.7, 0.3)$ as the default for arterial intersections
with asymmetric demand.

\begin{thebibliography}{9}
\bibitem{coffman1971system}
  E.~G. Coffman, M.~J. Elphick, and A.~Shoshani,
  ``System deadlocks,''
  \textit{ACM Computing Surveys}, vol.~3, no.~2, pp.~67--78, 1971.

\bibitem{tassiulas1992stability}
  L.~Tassiulas and A.~Ephremides,
  ``Stability properties of constrained queueing systems and scheduling policies
  for maximum throughput in multihop radio networks,''
  \textit{IEEE Transactions on Automatic Control}, vol.~37, no.~12,
  pp.~1936--1948, 1992.

\bibitem{demers1989analysis}
  A.~Demers, S.~Keshav, and S.~Shenker,
  ``Analysis and simulation of a fair queueing algorithm,''
  \textit{ACM SIGCOMM Computer Communication Review}, vol.~19, no.~4,
  pp.~1--12, 1989.
\end{thebibliography}

\end{document}